\section{从单轮到 Agent：为什么 agent benchmark 噪声更难}

\subsection{轨迹级方差（trajectory-level variance）}
在 agent benchmark 中，分数不仅受最终答案影响，还受中间决策路径影响。一个早期小错误会在后续工具调用中放大，形成高方差尾部事件（long-tail failure）。这使得 ``同一模型、同一任务'' 在多次运行中的结果分布更宽。

\begin{importantbox}{Agent eval 的本质变化}
从 ``静态样本打分'' 转向 ``动态轨迹打分''。评测对象不再是单输出，而是决策过程。
\end{importantbox}

\subsection{环境不确定性与评测一致性}
Agent 任务依赖外部环境：API 延迟、网页变化、工具可用性、文件状态都可能改变结果。因此课程强调要把 environment versioning 当成 benchmark 协议的一部分，必要时使用 sandbox replay 或 deterministic simulator。

\begin{knowledgebox}{Agent benchmark 的最小可复现要素}
\begin{itemize}
    \item environment snapshot 版本号；
    \item tool API mock/real 的开关与记录；
    \item timeout 与 retry 策略；
    \item trajectory log 与 error taxonomy。
\end{itemize}
\end{knowledgebox}

\subsection{评估指标需要从 ``accuracy'' 扩展到 ``reliability''}
Sida 讲中指出，对于 agent 系统，仅看 success rate 不够。至少还应报告 recovery rate、invalid action ratio、tool efficiency、time-to-success 等指标。这样才能区分 ``偶然成功'' 与 ``稳定可用''。

\begin{warningbox}{只看最终成功率会掩盖系统退化}
两个系统可能成功率接近，但一个需要大量无效工具调用和回退，另一个路径更短更稳。若只看 success rate，优化方向会偏离真实用户体验。
\end{warningbox}

\subsection{本章小结}
Agent benchmark 的难点不在 ``题更难''，而在 ``系统变量更多''。评测若不显式建模环境与轨迹噪声，就会把系统随机波动误判为模型能力变化。

